\documentclass[border=2pt]{standalone}
\usepackage{amsmath, amsthm, amsfonts}
\usepackage{tikz-cd, kotex}
\usepackage{quiver}
\usepackage{Operators}
\usepackage{palette}
\usepackage{diagram-style}
\begin{document}

%%FIG 복소곡면 n=2의 Hodge 다이아몬드: h^{p,q}를 마름모꼴로 배열, 수직축 좌우대칭이 Hodge 대칭 (Hodge_Theory-1.svg)
\begin{tikzpicture}
	\draw[densely dashed, black!50] (0,-0.7) -- (0,4.7);
	\node (a) at (0,4) {$h^{0,0}$};
	\node (b) at (-1.7,3) {$h^{1,0}$};
	\node (c) at (1.7,3) {$h^{0,1}$};
	\node (d) at (-3.4,2) {$h^{2,0}$};
	\node (e) at (0,2) {$h^{1,1}$};
	\node (f) at (3.4,2) {$h^{0,2}$};
	\node (g) at (-1.7,1) {$h^{2,1}$};
	\node (h) at (1.7,1) {$h^{1,2}$};
	\node (i) at (0,0) {$h^{2,2}$};
\end{tikzpicture}

\end{document}
